% Worked Examples: Concept 3.9.3 - System Design Considerations
\documentclass[11pt,a4paper]{article}
\usepackage{worked-examples-template}

\title{\textbf{Worked Examples}\\[0.5em]
\large Concept 3.9.3: System Design Considerations\\[0.3em]
\normalsize \coursesubtitle{} -- \coursetitle}
\author{Assoc.\ Prof.\ William Robertson \and Dr Sean McGowan}
\date{\today}

\begin{document}

\maketitle
\tableofcontents
\newpage

\section{Concept 3.9.3: System Design}

\subsection{Example 1: Stabilisability Check Using the Popov-Belevitch-Hautus Test}

The PBH test provides an algebraic method to determine if a system can be stabilised. This is fundamental for determining feasibility of control design.

\paragraph{Problem:}
Consider a system with state-space representation:
\[
A = \begin{bmatrix} 2 & 1 & 0 \\ 0 & -1 & 1 \\ 0 & 0 & 3 \end{bmatrix}, \quad
B = \begin{bmatrix} 0 \\ 1 \\ 0 \end{bmatrix}
\]

(a) Determine if the system is stabilisable using the Popov-Belevitch-Hautus (PBH) test.

(b) Identify which eigenvalues can be moved by state feedback.

(c) Explain why this system can or cannot be stabilised.

\paragraph{Solution:}

\textbf{Step 1: Find eigenvalues of $A$}

The eigenvalues are found from $\det(A - \lambda I) = 0$. Since $A$ is upper triangular, the eigenvalues are the diagonal elements:
\[
\lambda_1 = 2, \quad \lambda_2 = -1, \quad \lambda_3 = 3
\]

Eigenvalues in the right half-plane (unstable): $\lambda_1 = 2$ and $\lambda_3 = 3$.

\textbf{Step 2: Apply PBH test for stabilisability}

The PBH test states that a system is stabilisable if and only if:
\[
\operatorname{rank}([A - \lambda I \quad B]) = n
\]
for all eigenvalues $\lambda$ with $\operatorname{Re}(\lambda) \geq 0$ (right half-plane).

We need to check this for $\lambda_1 = 2$ and $\lambda_3 = 3$.

\textbf{For $\lambda_1 = 2$:}
\[
A - 2I = \begin{bmatrix} 0 & 1 & 0 \\ 0 & -3 & 1 \\ 0 & 0 & 1 \end{bmatrix}
\]
\[
[A - 2I \quad B] = \begin{bmatrix} 0 & 1 & 0 & 0 \\ 0 & -3 & 1 & 1 \\ 0 & 0 & 1 & 0 \end{bmatrix}
\]

Row reduce to find rank:
\[
\begin{bmatrix} 0 & 1 & 0 & 0 \\ 0 & -3 & 1 & 1 \\ 0 & 0 & 1 & 0 \end{bmatrix} \sim \begin{bmatrix} 0 & 1 & 0 & 0 \\ 0 & 0 & 1 & 1 \\ 0 & 0 & 1 & 0 \end{bmatrix} \sim \begin{bmatrix} 0 & 1 & 0 & 0 \\ 0 & 0 & 1 & 0 \\ 0 & 0 & 0 & 1 \end{bmatrix}
\]

Rank = 3 ✓ (satisfies PBH test for $\lambda = 2$)

\textbf{For $\lambda_3 = 3$:}
\[
A - 3I = \begin{bmatrix} -1 & 1 & 0 \\ 0 & -4 & 1 \\ 0 & 0 & 0 \end{bmatrix}
\]
\[
[A - 3I \quad B] = \begin{bmatrix} -1 & 1 & 0 & 0 \\ 0 & -4 & 1 & 1 \\ 0 & 0 & 0 & 0 \end{bmatrix}
\]

Row reduce:
\[
\begin{bmatrix} -1 & 1 & 0 & 0 \\ 0 & -4 & 1 & 1 \\ 0 & 0 & 0 & 0 \end{bmatrix}
\]

Rank = 2 ✗ (fails PBH test for $\lambda = 3$)

\textbf{Step 3: Conclusion}

The system is \textbf{NOT stabilisable} because the PBH test fails for the unstable eigenvalue $\lambda_3 = 3$.

Physically, this means:
\begin{itemize}
\item The eigenvalue $\lambda = 2$ is reachable and can be moved by state feedback
\item The eigenvalue $\lambda = 3$ is \emph{unreachable} from the input and cannot be affected by any control
\item Since $\lambda = 3$ is unstable and cannot be controlled, the system cannot be stabilised
\end{itemize}

\textbf{Key insights:}
\begin{itemize}
\item Stabilisability requires that all unstable modes be reachable from the input
\item A system can have unreachable modes, but only if they are stable
\item The PBH test provides a simple algebraic check without computing eigenvectors
\end{itemize}

\subsection{Example 2: Bode's Integral Formula and the Waterbed Effect}

Bode's Integral Formula reveals fundamental limitations: reducing sensitivity at one frequency inevitably increases it elsewhere.

\paragraph{Problem:}
Consider a feedback system with an unstable process $P(s) = 1/(s-1)$ and proportional controller $C(s) = k$. 

(a) For $k = 5$, calculate the sensitivity function $S(s)$ and loop transfer function $L(s)$.

(b) Evaluate $\log|S(\ii\omega)|$ at $\omega = 0.5$, $\omega = 2$, and $\omega = 10$ rad/s.

(c) Apply Bode's Integral Formula to find $\int_0^\infty \log|S(\ii\omega)|\,\mathrm{d}\omega$.

(d) Explain the waterbed effect for this system.

\paragraph{Solution:}

\textbf{Step 1: Loop transfer function and sensitivity}

Loop transfer function:
\[
L(s) = P(s)C(s) = \frac{k}{s-1} = \frac{5}{s-1}
\]

Sensitivity function:
\[
S(s) = \frac{1}{1 + L(s)} = \frac{1}{1 + 5/(s-1)} = \frac{s-1}{s-1+5} = \frac{s-1}{s+4}
\]

\textbf{Step 2: Evaluate $\log|S(\ii\omega)|$ at specific frequencies}

At $\omega = 0.5$ rad/s:
\[
S(\ii \cdot 0.5) = \frac{\ii \cdot 0.5 - 1}{\ii \cdot 0.5 + 4} = \frac{-1 + \ii \cdot 0.5}{4 + \ii \cdot 0.5}
\]
\[
|S(\ii \cdot 0.5)| = \frac{\sqrt{1 + 0.25}}{\sqrt{16 + 0.25}} = \frac{\sqrt{1.25}}{\sqrt{16.25}} = \frac{1.118}{4.031} = 0.277
\]
\[
\log|S(\ii \cdot 0.5)| = \log(0.277) = -1.284
\]

At $\omega = 2$ rad/s:
\[
|S(\ii \cdot 2)| = \frac{\sqrt{1 + 4}}{\sqrt{16 + 4}} = \frac{\sqrt{5}}{\sqrt{20}} = \frac{2.236}{4.472} = 0.500
\]
\[
\log|S(\ii \cdot 2)| = \log(0.5) = -0.693
\]

At $\omega = 10$ rad/s:
\[
|S(\ii \cdot 10)| = \frac{\sqrt{1 + 100}}{\sqrt{16 + 100}} = \frac{\sqrt{101}}{\sqrt{116}} = \frac{10.05}{10.77} = 0.933
\]
\[
\log|S(\ii \cdot 10)| = \log(0.933) = -0.069
\]

\textbf{Step 3: Apply Bode's Integral Formula}

Bode's Integral Formula states:
\[
\int_0^\infty \log|S(\ii\omega)|\,\mathrm{d}\omega = \pi \sum p_k
\]
where $p_k$ are the right half-plane poles of $L(s)$.

For our system, $L(s) = 5/(s-1)$ has one right half-plane pole at $p_1 = 1$.

Therefore:
\[
\int_0^\infty \log|S(\ii\omega)|\,\mathrm{d}\omega = \pi \cdot 1 = \pi \approx 3.14
\]

\textbf{Step 4: Interpretation - The waterbed effect}

\begin{itemize}
\item The integral is \emph{constant} and positive for unstable systems
\item At low frequencies ($\omega = 0.5$), $\log|S| < 0$ means $|S| < 1$ (good disturbance rejection)
\item At high frequencies ($\omega = 10$), $\log|S| > 0$ means $|S| > 1$ (amplification)
\item The negative area (good performance) at low frequencies must be balanced by positive area (poor performance) at higher frequencies
\item This is the "waterbed effect": pushing down sensitivity at one frequency causes it to pop up elsewhere
\end{itemize}

If we tried to increase $k$ to improve low-frequency disturbance rejection, we would:
\begin{itemize}
\item Reduce $|S|$ at low frequencies (more negative contribution to integral)
\item Force $|S|$ to increase more at mid-range frequencies (to maintain constant integral)
\item Create a larger sensitivity peak, reducing robustness
\end{itemize}

\textbf{Key insights:}
\begin{itemize}
\item Unstable systems have inherent performance limitations
\item The total "cost" of disturbance rejection is fixed by the unstable pole location
\item Control design is always a compromise between performance at different frequencies
\item Stable systems ($\sum p_k = 0$) have zero integral, allowing better performance flexibility
\end{itemize}
\end{document}
